\section{A double curve of genus two}\label{sec:double}

In genus two the first difference between the two circle actions can
be calculated inside the complex cycle algebra. The identity also
follows under numerical reconstruction in \cite[\S6]{PaperIII}.

\begin{proposition}\label{prop:double}
In $H$ one has $a_{2,2}=x_{2,2}-2x_{4,1}$.
\end{proposition}
\begin{proof}
Let $C$ be a smooth genus-two curve with hyperelliptic involution
$\iota$. Choose distinct non-Weierstrass points $p,q$ in distinct
orbits. Choose $L$ with $L^2=\mathcal O(p+q)$, possible by
surjectivity of multiplication by two on the Jacobian. Put
$M=K_C L^{-1}$. The divisors of chosen sections $s_L,s_M$ of
$L^2,M^2$ are $p+q,\iota p+\iota q$, respectively. Indeed
$p+\iota p\sim q+\iota q\sim K_C$. Each section spans its space
of global sections by Riemann--Roch, since neither divisor is
canonical. A section of $L$ would give a square section with a
double zero, which is impossible. The same holds for $M$. Also
$ML^{-1}=K_C(-p-q)$ is a nontrivial degree-zero line bundle.
By the deformation of balanced pairs in \S\ref{sec:local} and the
scalar construction \cite[\S4.3]{Pardon}, both classes can be
computed on this splitting.

In $Y=\operatorname{Tot}_C(L\oplus M)$, every irreducible compact
curve projects finitely onto $C$. A degree-two cycle has zero trace
in both fiber coordinates $u,v$, because $H^0(L)=H^0(M)=0$.
Its two fiber points, with multiplicity, thus have the form
$\pm(u,v)$. The three quadratic invariants give global sections

\[
        u^2=a s_L,\qquad v^2=b s_M,\qquad uv=d\in H^0(K_C).
\]

They extend over ramification by symmetric evaluation on the finite
maps from the normalizations of the components. Since
$d^2=ab\,s_Ls_M$ and the four zeros of $s_Ls_M$ are simple,
$ab=0$ and $d=0$. Conversely each point of the two coordinate
axes gives a cycle. On the first axis it is
$D_a=\{u^2=a s_L,\ v=0\}$, and there is an exchanged formula
on the second. The origins are the same cycle $2[C]$. These
parametrizations are continuous, including at zero, by the finite
flat families of spectral divisors. Their inverse is continuous by
fiber evaluation away from the zeros. This identifies the underlying
cycle space equivariantly with the union of the two complex lines.
For nonzero amplitude each cover is smooth and connected, with
two simple branch points. Riemann--Hurwitz gives genus four. The
two independent scalings of $L,M$ act on $a,b$ with characters
$(2,0),(0,2)$.

\smallskip\noindent\emph{The two Borel classes.}
Set $B=\mathbb P^1\times\mathbb P^1$ and let $h_1,h_2$ be the
positive hyperplane classes. Twist $L,M$ by the hyperplane lines
from the first and second factors, respectively. The relative
degree-two cycle space $W$ is the union of the total spaces
$P_+=\operatorname{Tot}\mathcal O(2,0)$ and
$P_-=\operatorname{Tot}\mathcal O(0,2)$, glued along their zero
sections $B$. Write $U_+,U_-$ for the complex Thom classes of
these two line bundles. Use the two base inclusions

\[
            i_s(z)=(z,z),\qquad i_a(z)=(z,\bar z),
\]

orienting the source by its complex orientation in both cases.
The first restricts the target family to the scalar Borel family.
On the second, the second hyperplane line identifies real analytically
and fiberwise complex linearly with $\mathcal O(-1)$ by the
Hermitian metric. Thus its restriction is the antidiagonal Borel
family. A real analytic base map and real analytic bundle
identification suffice here. On real simplices they complexify to
the family germs of \cite[Definition~3.1]{Pardon}. The Gysin map of either
embedding preserves the universal image. One can use a product
triangulation subdivided along the diagonal, and transport it by
conjugation of the second factor for the second graph. Both give
the canonical concordance class of immersed real analytic
triangulations of $B$. In the shifted compactly supported
description the pushforward is the oriented normal Thom map, as
in Pardon's diagonal computation~\cite[\S3.1, Lemma 3.5 and proof of Lemma 4.13]{Pardon}.

At Borel stage one the scalar cut is $(a+b)^*\mathrm{Th}/2$,
where the sum is in $\mathcal O(2)$. For the antidiagonal cut
use $a+j(b)$, where the Hermitian conjugate-linear identification
$j:\mathcal O(-2)\to\mathcal O(2)$ reverses the real fiber
orientation. Again divide by two. Each cut vanishes exactly on the
common zero section. These choices represent the published
normalizations. To check this even for the continuous test $j$,
use the standard Hermitian metrics at every finite Borel stage.
The tests restrict along hyperplanes, and each ordinary normalized
Euler class is $h$. Append a defining tuple from \cite[\S2]{PaperIII} to the
one-test tuple and append this test to the defining tuple. The two
resulting compactly supported Thom products coincide at the common
stage by naturality and the projection formula of the stabilization
lemma. The same argument identifies the scalar test with that of \cite[Definition~4.7]{Pardon}.

The Poincaré duals of the two base graphs are $h_1+h_2$ and
$-h_1+h_2$, respectively. For the second sign, integrate $h_1$
and $h_2$ on the graph to obtain $1,-1$. Naturality of the
normal Thom map under restriction to either branch gives the
following table in $H_c^4(P_\pm;\mathbb Q)$.
\[
\begin{array}{lcc}
 & P_+ & P_- \\ \hline
\text{scalar} & (h_1+h_2)U_+/2 &(h_1+h_2)U_-/2\\
\text{antidiagonal} & (-h_1+h_2)U_+/2 &(h_1-h_2)U_-/2\\
\text{difference} &-h_1U_+ &-h_2U_-
\end{array}
\]
These restrictions determine the class on $W$. For the proper
map $\nu:P_+\sqcup P_-\to W$ the difference of restrictions
to $B$ gives the exact sequence
$0\to\mathbb Q_W\to\nu_*\mathbb Q_{P_+\sqcup P_-}
\to\mathbb Q_B\to0$. There are no higher direct images since the
fibers are finite. Now $H^3(B)=0$, so restriction of $H_c^4(W)$
to the two branches is injective.

For the positive class $(h_1U_+,0)$, fix one point on the first
base factor, trivialize its hyperplane line there, and fix a nonzero
amplitude $a_*$. Over the remaining base $\mathbb P^1$, take
the complex point class of $a_*$ in the amplitude line and extend
by zero from a disk avoiding zero. Pushing forward from the base
slice gives precisely $(h_1U_+,0)$ by the Thom isomorphism. Its
support misses the other branch. The exchanged construction
represents $(0,h_2U_-)$. It remains to compute the universal
image of the class on each base slice.

\smallskip\noindent\emph{The class of a branch.}
We treat the first branch. The family over the remaining
$B'=\mathbb P^1$ has first summand $L$ and second summand
$M\otimes\mathcal O_{B'}(1)$. Write $S=\operatorname{Tot}_C L$.
Deform the threefolds as affine line bundles over $S$, retaining
the projection to $S$. Choose independent sections
$\sigma_0,\sigma_1$ of $\mathcal O_{B'}(1)$. Let $z_0$ be
the simple zero of $\sigma_0$. Choose a nonzero class

\[
     c\in H^1(C,ML^{-1})\simeq\mathbb C .
\]

Cover $C$ by a disk $V$ avoiding the zeros of $s_L$, and the
complement $U$ of its center. These two Stein sets and their
intersection compute this cohomology. Represent $c$ on $U\cap V$.
For $t\in[0,1]$, add the following translation to the usual
transition of the second fiber coordinate, written intrinsically
in the common line over the overlap. The added term is

\[
     t\,c\,u\,
          \big[\sigma_0+(u^2/s_L-a_*)\sigma_1\big].
\]

This is holomorphic on the whole overlap, including $u=0$.
The two-set gluing with inverse translation defines a holomorphic
family, also for complex $t$. The relative tangent determinant
in each threefold is the pullback of $TC\otimes L\otimes M$, up
to a constant line on the fiber, so the cycles stay semi-Fano.

Fix $a$ in a small disk about $a_*$ avoiding zero. Any cycle
of degree two with this amplitude projects onto the smooth
spectral divisor $D_a\subset S$ with degree one. Indeed its
projected cycle is determined by $u^2=a s_L$, and no component
can lie in an affine line fiber. The finite degree-one map to the
smooth divisor is an isomorphism on the normalization and the lift
is a section of the affine line bundle over $D_a$. Conversely
every such section gives the embedded cycle. Write
$\pi:D_a\to C$. The identity
$\pi_*\mathcal O_{D_a}=\mathcal O_C\oplus L^{-1}$ gives

\[
   H^0(D_a,\pi^*M)=0,\qquad
   H^1(D_a,\pi^*M)=H^1(C,ML^{-1}) .
\]

The lifting obstruction is
$t[\sigma_0(z)+(a-a_*)\sigma_1(z)] [c\,u]$.
Here multiplication by $u$ identifies $c$ with its nonzero
class in the odd summand of the displayed direct image. Hence a
lift exists precisely when this expression vanishes. When it does,
the translations on the divisor are themselves zero, and the zero
section in the two local charts gives the unique lift.

Pull back the Thom class of $a-a_*$ on this amplitude disk. It
defines a supported class throughout the prism $B'\times[0,1]$.
Its zero set consists of the section $D_{a_*}$ over all of $B'$
at $t=0$, together with $D_{a_*}$ over $z_0$ for every $t$.
This set is compact. The amplitude is continuous by fiber
evaluation in the projection to $S$. On the zero-obstruction
locus the explicit local zero graphs vary continuously by their
spectral equations. Thus the parametrization above is also a
homeomorphism over the disk. Extension by zero and the oriented
prism relation identify the endpoint insertions. At $t=1$, the
relative cycle space near the support is the graph

\[
                  a=a_*-\sigma_0(z)/\sigma_1(z).
\]

The function on the right has nonzero derivative at $z_0$.
The insertion has local degree $+1$ on this graph. Projection
of the graph to the base is a complex orientation preserving local
isomorphism, so the Gysin image from the fiber over $z_0$ of the
unit on the isolated curve is exactly this class.

For regularity in that fiber, the normal sequence of $D=D_{a_*}$
has subbundle $\pi^*M$ and quotient $\pi^*L^2$. The latter is
the normal line to $D$ in $S$, by its divisor equation.
Its space of sections is generated by the amplitude variation
$da=1$, since $H^0(L^2)=\mathbb C s_L$ and $H^0(L)=0$.
The normal connecting map sends this variation to
$\pm\sigma_1(z_0)[c\,u]$, by differentiating the transition
on $D_a$ at $a=a_*$. It is an isomorphism. Also
$H^1(D,\pi^*L^2)=H^1(L^2)\oplus H^1(L)=0$. It follows that
$H^0(N_{D/Y_{z_0,1}})=H^1(N_{D/Y_{z_0,1}})=0$.

Finally deform the ambient germ along this smooth curve to the
total space of its normal bundle $N$. One model is the blow-up
of $Y_{z_0,1}\times\mathbb C$ along $D\times0$, with the
strict transform of the central fiber removed. The copy of $D$
persists and is isolated and regular in every fiber, including
the normal bundle fiber. Near a smooth embedded curve of multiplicity
one the relative cycles are embedded deformations, and vanishing
of both normal cohomology groups makes the relative space locally
the base. Thus the section over a closed parameter interval is
compact and open in the relative cycle space there, and its unit
continues to the zero section of $N$. Every degree-one compact
cycle in $N$ is a holomorphic section, so its degree-one cycle
space is the single point $H^0(D,N)=0$. In \cite[Definition~4.7]{Pardon} one can take Borel stage zero and the empty tuple,
since the whole degree-one locus is compact and its class is the unit.
To identify the indices, adjunction gives
$\deg N=2g(D)-2=6$, since the fiber threefold has trivial
canonical bundle. Thus $c_1(T\operatorname{Tot}N)[D]=0$.
Pardon's scalar construction applies to any rank-two holomorphic
bundle on a smooth proper curve of this genus and Chern number.
This includes the possibly nonsplit $N$ obtained from the normal
sequence. With $m=1$ the empty tuple at stage zero therefore gives
$+x_{4,1}$, with the sign fixed by the local degree on the graph~\cite[Remark 4.1 and Definition 4.7]{Pardon}.

Exchanging $L$ and $M$ preserves the assumptions, including
nontriviality of the degree-zero ratio, so the second positive
branch class has the same image. Substitution in the last row of
the table gives $-2x_{4,1}$ for the difference.
\end{proof}
